\section{引言}
磁驱动能够在细长机器人末端不安装电机的情况下施加力和力矩，但用于血管介入时，形状、材料姿态、周围流体与执行器几何共同决定运动，因而对模型提出了较高要求。本文的长期应用目标是空间站微重力条件下的三维柔性血管内机器人。重力主导载荷减弱，并不意味着血流作用、壁面相互作用、磁工作空间约束或转动惯性消失。本文研究支撑这一目标所需的数学与计算接口；证据来自仿真和局部分析，尚不构成生理或临床验证。

只有先说明保持的是哪一动力学对象的结构，保结构更新才具有明确意义。正则 Hamiltonian 降阶在保守分支上给出一种结构，非保守载荷通常破坏其保持性。时间相关 Birkhoff 形式可以局部表示光滑受载向量场，但近似输运形式可能表示的是邻近向量场。同样，学习得到的标量能量定义了保守预测器，而磁执行器只能施加受限的动量广义力。混淆这些对象，可能在结构缺陷很小的同时产生不准确的物理轨迹。对于具有位形相关质量与空间变化线圈场的三维杆，这一区别尤其重要。

本文以包含两向弯曲和扭转的几何 Kirchhoff 降阶模型研究上述接口，并通过空间 Cosserat 模型考察环境转移。降阶模型保留12个配置变量及质量矩阵的全部导数。完整 Biot--Savart 磁能给出电流到广义力的映射，随后在明确的光滑分支上评估局部一形式输运、Moser 图和端点生成函数。Hamiltonian 学习与共振诊断共用理想偏置保守降阶模型，完整线圈闭环则单独标明。因此，本文给出的是具有明确适用范围的模型接口连接，而不是将所有方法合并为一项实验。

本文成果分为三部分。首先，建立从变惯性杆向量场到输运二形式及相容生成更新的可计算局部连接，分别保留表示、坐标图、积分与求解器误差。推导明确端点横截性、初始规范及固定坐标递推，并在自治保守分支恢复非可分四阶中心方法。这是在既有几何力学框架中的应用与相容条件展开，不将 Darboux 或生成函数理论本身作为本文首创。其次，把磁势、虚功、反馈求导和受限分配连接到同一几何降阶模型，并量化局部仿射场在所测线圈几何中的失效。最后，报告学习、共振与空间杆转移的边界：有限样本上的场预测准确不保证动量曲率为正；形式共振 twist 非零不构成 KAM 证书；端口执行满足约束也不保证跟踪改善。

\section{相关工作及研究记录的适用范围}
几何精确杆模型区分中心线变形与材料标架转动。\citet{Antman2005,Simo1985} 给出其变分与计算基础，\citet{Gazzola2018} 研究柔性细丝的正、逆力学。这些模型支持在完整 Cosserat 计算中保留剪切与伸长，同时在低维解析模型中明确施加 Kirchhoff 约束。\citet{Abbott2020} 综述磁驱动方法及其工作空间依赖性；铁磁连续体机器人展示了相应物理平台 \citep{Kim2019}，受限管腔中的磁软体机器人仿真与导航则为本文参数及实现来源提供参考 \citep{Tong2025}。本文线圈阵列与残余加速度仍是基准设定，而非硬件测量或微重力生理数据。

生成函数可用于构造正则差分格式 \citep{Feng1989}，辛算法的系统理论见 \citet{FengQin2010,Hairer2006}。Birkhoff 格式将结构问题推广到非正则二形式 \citep{Su2007}。对于耗散物理动力学，二形式的时间相容性不可缺少：空间闭性本身不能证明其表示给定向量场。本文通过辅助 Moser 流计算局部坐标图，并保留图的时间导数。原机器人精确流只用于局部存在性论证，不作为可用数值输入。本文四阶格式是适用于非可分 Hamiltonian 的中心生成方法，而非可分系统的漂移--踢步组合。

Hamiltonian 神经网络通过能量的导数定义向量场 \citep{Greydanus2019}。从学习进入 KAM 的论证需要比有限测试损失强得多的范数控制与假设 \citep{Chen2022}。经典与有限可微持久性结果 \citep{Poschel2001,Bounemoura2020}，以及辛算法的离散结果 \citep{Shang1999}，明确了这些证明责任。本文圆截面杆存在四对相等弯曲频率，因此不能直接复用非共振两模态证书。本文将既有有限特征及统计推论保留为条件性结果，报告已计算的共振项，而不宣称得到新的不变环面。

本文证据截止于 Stage~5 分支的仓库提交 \texttt{1c73a58}。原始1025稿件提供物理与数学意图，Stage~4审查指出必要修正，Stage~5提供实现及已有验证。本文不使用 Stage~6 或后续衔接研究的新结果，不新增科学仿真或训练。已有输出按其记录的验证状态选用，被否定的导数计算及失败分配保留独立身份。随附清单把公式、数值主张及图中序列对应到精确源文件与字段。这一范围限制很重要：仓库中历史两模态及平面实验的存在，不能证明新的三维结论。
